% ─── Slide 12: Post-ADC LoRA ─────────────────────────────────────────────────
\begin{frame}[shrink=10]{Method II: Post-ADC Residual LoRA}

  \textbf{Key idea.} Freeze everything PTQ produced. Add a tiny low-rank branch that sits \alert{in parallel} with the ADC-quantised projection, reads the FP input, and adds to the ADC output:
  \[
    y \;=\; \underbrace{\mathrm{ADC}_{\mathrm{frozen}}\!\left(\hat x\right)}_{\text{quantised path}} \;+\; \underbrace{\tfrac{\alpha_{\scriptscriptstyle\mathrm{LoRA}}}{r}\, B\bigl(A\bigl(x_{\mathrm{fp}}\bigr)\bigr)}_{\text{residual correction, rank }r}
  \]

  \vspace{-0.2em}
  \begin{columns}[T]
    % ── LoRA-paper-style diagram ─────────────────────────────────
    \begin{column}{0.5\textwidth}
      \begin{center}
      \resizebox{0.98\linewidth}{!}{%
      \begin{tikzpicture}[
          every node/.style={font=\footnotesize},
          box/.style={draw, rounded corners=2pt, minimum width=1.6cm, minimum height=0.6cm, align=center},
          frozen/.style={fill=mygray!15, draw=mygray!60},
          traina/.style={fill=myorange!25, draw=myorange!80, thick},
          sum/.style={draw, circle, minimum size=0.55cm, inner sep=0pt},
          lab/.style={font=\tiny, text=mygray}
        ]

        \node (x) at (0, 3.5) {$x_{\mathrm{fp}}$};

        % Left branch: frozen ADC path
        \node[box, frozen] (qx) at (-1.7, 2.5) {quantise $\hat x$};
        \node[box, frozen] (w)  at (-1.7, 1.5) {$\hat W$ \small(frozen)};
        \node[box, frozen] (adc) at (-1.7, 0.5) {ADC (frozen)};

        \draw[-{Latex}] (x) -| (qx);
        \draw[-{Latex}] (qx) -- (w);
        \draw[-{Latex}] (w) -- (adc);

        % Right branch: LoRA
        \node[box, traina] (a) at (1.7, 2.0) {$A\in\mathbb R^{r\times d}$\\ \scriptsize trainable};
        \node[box, traina] (b) at (1.7, 0.9) {$B\in\mathbb R^{d\times r}$\\ \scriptsize trainable};
        \draw[-{Latex}] (x) -| (a);
        \draw[-{Latex}] (a) -- (b);

        % Sum
        \node[sum] (plus) at (0, -0.3) {$+$};
        \draw[-{Latex}] (adc.south) |- (plus);
        \draw[-{Latex}] (b.south) |- (plus);

        \node[below=0.15cm of plus] (y) {$y$};
        \draw[-{Latex}] (plus) -- (y);

        \node[lab, left=0.05cm of adc.west] {analog path};
        \node[lab, right=0.05cm of b.east] {digital rank-$r$};
      \end{tikzpicture}}
      \end{center}
      {\scriptsize \emph{Analogous to the original LoRA figure} \textit{(Hu et al.\ 2021)}, \emph{but added \textbf{after} the ADC, not to pre-trained weights.}}
    \end{column}

    % ── right column: loss and trainable params ─────────────────
    \begin{column}{0.48\textwidth}
      \footnotesize
      \textbf{Loss: distill from FP teacher.}
      \[
        \mathcal L_{\mathrm{LoRA}} \;=\; \mathrm{CE}(\hat p, y) \;+\; \lambda\,\mathrm{KL}\!\left(\hat p\,\Vert\, p_{\mathrm{FP}}\right)
      \]
      \vspace{-0.3em}
      \scriptsize $\lambda{=}1$, $\alpha_{\scriptscriptstyle\mathrm{LoRA}}/r$ scaling as in the LoRA paper.

      \vspace{0.4em}
      \pause
      \textbf{What is trained:}
      \begin{itemize}\scriptsize
        \item Only $A, B$ for each of 7 projections ($q,k,v,o,\text{gate},\text{up},\text{down}$).
        \item Everything else (weights, rotations, diagonals, ADC params) is \alert{frozen}.
        \item FP teacher runs once to produce $p_{\mathrm{FP}}$; no gradient through it.
      \end{itemize}

      \pause
      \vspace{0.3em}
      \begin{exampleblock}{Why post-ADC works}
        \scriptsize
        Pre-ADC LoRA \emph{fails} — its correction gets quantised by the fixed $\delta$. Post-ADC, the correction is added in the digital domain with FP precision.
      \end{exampleblock}
    \end{column}
  \end{columns}
\end{frame}
